%============================================================
%  ch12_quantum_states.tex
%  Chapter 12: Quantum States as Tensor Formats
%
%  Tensor Formats in Banach Spaces
%  Antonio Falcó, CEU Universities
%
%  Based on: Acevedo--Falcó (2026), FH2012, FHN2019
%  Estimated length: ~25 pp.
%============================================================

\chapter{Quantum States as Tensor Formats}
\label{chap:quantum_states}

The theory developed in Parts~I--IV connects naturally to quantum
mechanics through a single observation: for a non-zero vector $v$ in a
Hilbert tensor space $H = \bigotimes_{j=1}^d H_j$, the minimal subspace
$\Umin{\alpha}{v}$ coincides with the support of the reduced density
operator of the quantum state associated to $v$. This chapter makes
this connection precise and develops the geometric theory of quantum
states from the perspective of tensor formats and minimal subspaces.

\begin{remark}[Notation convention for Part~\ref{part:quantum}]
\label{rem:qs:notation}
In Parts~\ref{part:foundations}--\ref{part:variational}, tensors in
$\VDnorm$ are written in boldface ($\bv, \bu, \bw$) to distinguish
them from their component vectors ($v_\alpha \in V_\alpha$). In this
part, the primary objects are density operators $\rho \in \cTsa{1}(H)$
and Hilbert-space vectors $v \in H$; we follow the standard convention
of the operator-theory literature and use lightface notation throughout.
When a pure state $\rho_v = v \otimes J(v)/\|v\|^2$ is considered as
a tensor in $H \cong \bigotimes_j H_j$ and the tensor structure is the
explicit focus, the boldface notation $\bv$ is used.
\end{remark}

The state space of a quantum system on $H$ is the set $\Dstate{H}$
of positive trace-class operators of unit trace (see
Appendix~\ref{app:quantum} for the relevant operator-theoretic background).
Each fixed-rank stratum $\Drank{r}{H}$ carries the structure of a smooth
Banach manifold, fibred over the Grassmannian $\Grass{r}{H}$ — precisely
the structure of the Tucker manifold of Chapter~\ref{chap:tucker_manifold}
specialised to the trace-class Banach space $\cTsa{1}(H)$.

Section~\ref{sec:qs:pure} establishes the fundamental connection between
minimal subspaces and partial traces for pure states.
Section~\ref{sec:qs:mixed} extends to mixed states and identifies
the rank stratification of $\Dstate{H}$ with the Tucker rank stratification.
Section~\ref{sec:qs:entanglement} connects the Tucker rank vector with
the entanglement structure of multipartite systems.
Section~\ref{sec:qs:manifold} develops the Banach manifold structure of
$\Drank{r}{H}$ as a special case of the Tucker manifold theory.
Section~\ref{sec:qs:approx} applies the best-approximation theory of
Chapter~\ref{chap:minsubspaces} to quantum state compression.
Section~\ref{sec:qs:history} contains historical notes.

Throughout this chapter, $H = \bigotimes_{j=1}^d H_j$ is a Hilbert
tensor product of separable complex Hilbert spaces $H_j$,
and $\cTsa{1}(H)$ is the real Banach space of self-adjoint trace-class
operators on $H$, equipped with the trace norm $\|\cdot\|_1$
(see Appendix~\ref{app:quantum}).

%------------------------------------------------------------
\section{Pure states and minimal subspaces}
\label{sec:qs:pure}
%------------------------------------------------------------

\subsection*{The pure state associated to a vector}

For a non-zero vector $v \in H$, the \emph{pure state associated to $v$} is
\begin{equation}
\label{eq:qs:pure_state}
    \rho_v \coloneqq \frac{v \otimes J(v)}{\|v\|^2}
    \in \Dstate{H},
\end{equation}
where $J: H \to H^*$ is the Riesz isomorphism and $v \otimes J(v)$
is the rank-one operator $w \mapsto \langle v, w\rangle v / \|v\|^2$.
Two vectors $v$ and $\lambda v$ ($\lambda \in \CC^*$) give the same
pure state. The map $v \mapsto \rho_v$ thus induces a bijection between
the projective space $\mathbb{P}(H)$ and the set $\Drank{1}{H}$ of
pure states.

\subsection*{The fundamental connection}

The following theorem is the cornerstone of this chapter: it identifies
the minimal subspaces of $v$ (in the sense of
Definition~\ref{def:minimal_subspace}) with the supports of the
partial traces of $\rho_v$.

\begin{theorem}[Minimal subspaces and partial traces]
\label{thm:qs:minimal_partial_trace}
Let $v \in H = \bigotimes_{j=1}^d H_j$ with $\|v\|=1$ and let
$\rho_v = v \otimes J(v)$ be the associated pure state. For each
$\alpha \in D = \{1, \ldots, d\}$, let
$\rho_\alpha \coloneqq \TrS{\alpha^c}(\rho_v) \in \cTsa{1}(H_\alpha)$
be the partial trace over all factors except $\alpha$. Then:
\begin{enumerate}
\item[\textup{(a)}] $\rho_\alpha$ is a positive trace-class operator
    with $\Tr(\rho_\alpha) = 1$ and
    \begin{equation}
    \label{eq:qs:supp_min}
        \supp(\rho_\alpha) = \Umin{\alpha}{v}.
    \end{equation}
\item[\textup{(b)}] $\mathrm{rank}(\rho_\alpha) = r_\alpha(v)$,
    the Tucker $\alpha$-rank of $v$.
\item[\textup{(c)}] The non-zero eigenvalues of $\rho_\alpha$ are
    $\lambda_k^{(\alpha)}(v) > 0$, $k = 1, \ldots, r_\alpha(v)$,
    with $\sum_k \lambda_k^{(\alpha)}(v) = 1$.
    For $d=2$, $\lambda_k^{(1)}(v) = \lambda_k^{(2)}(v) = s_k(v)^2$
    where $s_k(v)$ are the Schmidt coefficients of $v$.
\end{enumerate}
\end{theorem}

\begin{proof}
By the definition of $\Umin{\alpha}{v}$ (Definition~\ref{def:minimal_subspace}),
$\Umin{\alpha}{v}$ is the smallest subspace $U \subset H_\alpha$ such
that $v \in U \otimes_H H_{\alpha^c}$. This is exactly the range of the
map $T: H_\alpha^* \to H_{\alpha^c}$ defined by $T(\varphi) = (\varphi
\otimes \Id{\alpha^c})(v)$, hence $\Umin{\alpha}{v} =
\overline{\mathrm{Im}(T^*)}$ where $T^*: H_{\alpha^c}^* \to H_\alpha$
is the adjoint.

On the other hand, $\rho_\alpha = \TrS{\alpha^c}(\rho_v)$ is the
operator on $H_\alpha$ characterised by
$\langle u, \rho_\alpha w\rangle_{H_\alpha} = \langle u \otimes
\Id{\alpha^c}, \rho_v (w \otimes \Id{\alpha^c})\rangle$ for all
$u, w \in H_\alpha$. A direct computation using
$\rho_v = v \otimes J(v)$ gives
\begin{equation*}
    \rho_\alpha = \TrS{\alpha^c}(v \otimes J(v))
    = T T^*: H_\alpha \to H_\alpha
\end{equation*}
(the Gram operator of $T$ composed with its adjoint). Hence
$\mathrm{Im}(\rho_\alpha) = \mathrm{Im}(TT^*) = \overline{\mathrm{Im}(T^*)}
= \Umin{\alpha}{v}$, proving~\eqref{eq:qs:supp_min}.
Part~(b) follows since $\mathrm{rank}(TT^*) = \mathrm{rank}(T) = r_\alpha(v)$.
Part~(c) follows from the spectral theorem for compact operators
(Theorem~\ref{thm:app:spectral}): the eigenvalues of $\rho_\alpha = TT^*$
are the squared singular values of $T$, and $\Tr(\rho_\alpha) = 1$
gives $\sum_k \lambda_k^{(\alpha)} = 1$.
\end{proof}

\begin{remark}[The $d=2$ case: Schmidt decomposition]
\label{rem:qs:schmidt}
For $d=2$, Theorem~\ref{thm:qs:minimal_partial_trace} reduces to the
classical Schmidt decomposition. Every $v \in H_1 \otimes_H H_2$ can
be written uniquely (up to phases) as
\begin{equation*}
    v = \sum_{k=1}^r s_k\, e_k \otimes f_k,
    \quad s_1 \geq s_2 \geq \cdots \geq s_r > 0,
\end{equation*}
with $\{e_k\}_{k=1}^r \subset H_1$ and $\{f_k\}_{k=1}^r \subset H_2$
orthonormal, $r = r_1(v) = r_2(v)$. Then
$\Umin{1}{v} = \spann\{e_k\}$ and $\Umin{2}{v} = \spann\{f_k\}$,
and $\rho_1 = \sum_k s_k^2\, e_k \otimes J(e_k)$,
$\rho_2 = \sum_k s_k^2\, f_k \otimes J(f_k)$.
\end{remark}

\begin{remark}[Tucker format as the natural language for multipartite states]
\label{rem:qs:tucker_natural}
Theorem~\ref{thm:qs:minimal_partial_trace} shows that for a pure state
$\rho_v$, the Tucker rank vector $(r_\alpha(v))_{\alpha \in D}$ is
an intrinsic property of the state: it is determined by the ranks of
the reduced density operators $\{\rho_\alpha\}$, which are physical
observables. The Tucker manifold $\Tucker{\fr}{H}$ is therefore the
natural variational class for quantum systems with prescribed
entanglement structure $(r_\alpha)_{\alpha \in D}$.
\end{remark}

\begin{figure}[ht]
\centering
\begin{tikzpicture}[>=Latex, node distance=1.8cm,
  state/.style={draw, rounded corners=4pt, minimum width=2.4cm,
                minimum height=0.75cm, align=center, font=\small},
  obs/.style={draw, dashed, rounded corners=4pt, minimum width=2.4cm,
              minimum height=0.75cm, align=center, font=\small, fill=orange!8},
  arr/.style={->, thick},
  label/.style={font=\scriptsize, draw=none}
]
% Central: pure state
\node[state, fill=blue!8] (v)
  {$v \in \Tucker{\fr}{H}$\\[1pt]\scriptsize pure state};

% Tucker rank annotation
\node[label, above=0.4cm of v]
  {Tucker rank $\fr = (r_\alpha)_{\alpha \in D}$};

% Partial traces
\node[obs, left=3.0cm of v] (r1)
  {$\rho_1 = \TrS{1^c}(\rho_v)$\\[1pt]\scriptsize red.\ density op.\ on $H_1$};
\node[obs, right=3.0cm of v] (r2)
  {$\rho_2 = \TrS{2^c}(\rho_v)$\\[1pt]\scriptsize red.\ density op.\ on $H_2$};
\node[obs, below=1.6cm of v] (ralpha)
  {$\rho_\alpha = \TrS{\alpha^c}(\rho_v)$\\[1pt]\scriptsize
   for each $\alpha \in D$};

% Minimal subspaces
\node[state, fill=green!8, left=3.0cm of r1, yshift=0cm] (U1)
  {$\Umin{1}{v} = \supp(\rho_1)$\\[1pt]\scriptsize $\dim = r_1$};
\node[state, fill=green!8, right=3.0cm of r2, yshift=0cm] (U2)
  {$\Umin{2}{v} = \supp(\rho_2)$\\[1pt]\scriptsize $\dim = r_2$};
\node[state, fill=green!8, below=1.4cm of ralpha] (Ualpha)
  {$\Umin{\alpha}{v} = \supp(\rho_\alpha)$\\[1pt]\scriptsize $\dim = r_\alpha$};

% Arrows: pure state -> partial traces
\draw[arr, blue] (v) -- node[above, font=\scriptsize]{$\TrS{1^c}$} (r1);
\draw[arr, blue] (v) -- node[above, font=\scriptsize]{$\TrS{2^c}$} (r2);
\draw[arr, blue] (v) -- node[right, font=\scriptsize]{$\TrS{\alpha^c}$} (ralpha);

% Arrows: partial traces -> minimal subspaces
\draw[arr, orange!70!black] (r1)
  -- node[above, font=\scriptsize]{$\supp(\cdot)$} (U1);
\draw[arr, orange!70!black] (r2)
  -- node[above, font=\scriptsize]{$\supp(\cdot)$} (U2);
\draw[arr, orange!70!black] (ralpha)
  -- node[right, font=\scriptsize]{$\supp(\cdot)$} (Ualpha);

% Curved arrow: direct connection
\draw[arr, dashed, green!60!black] (v.south west)
  to[bend right=30]
  node[left, font=\scriptsize, green!50!black]{Def.~of $\Umin{\alpha}{\cdot}$}
  (U1.north east);

% Tucker manifold annotation at bottom
\node[label, below=0.3cm of Ualpha, align=center]
  {$v \in \bigotimes_\alpha \Umin{\alpha}{v}$,\quad
   $r_\alpha(v) = \mathrm{rank}(\rho_\alpha)$};
\end{tikzpicture}
\caption{The correspondence between minimal subspaces and partial traces
(Theorem~\ref{thm:qs:minimal_partial_trace}), illustrated for
$v \in \Tucker{\fr}{H} \subset H = \bigotimes_{j \in D} H_j$.
Blue arrows: formation of reduced density operators by partial tracing.
Orange arrows: extraction of minimal subspaces as supports of the
reduced operators. The dashed green arrow shows the direct definition
of $\Umin{\alpha}{v}$ as the smallest $U \subset H_\alpha$ with
$v \in U \otimes_H H_{\alpha^c}$.
The Tucker $\alpha$-rank $r_\alpha(v) = \dim\Umin{\alpha}{v}$ equals
the rank of the corresponding reduced density operator.}
\label{fig:minimal_partial_trace}
\end{figure}

%------------------------------------------------------------
\section{Mixed states and the rank stratification}
\label{sec:qs:mixed}
%------------------------------------------------------------

\subsection*{Hilbert--Schmidt operators as tensors}

By Theorem~\ref{thm:app:hs_tensor}, $\cT{2}(H) \cong H \otimes_H H^*$
isometrically as Hilbert spaces. Under this identification, a
Hilbert--Schmidt operator $A \in \cT{2}(H)$ of rank $r$ corresponds
to a tensor in $H \otimes_H H^*$ of Tucker rank $(r, r)$ with respect
to the bipartition $\{H, H^*\}$. The SVD of $A$,
\begin{equation*}
    A = \sum_{k=1}^r s_k(A)\, v_k \otimes J(u_k),
\end{equation*}
is the Tucker decomposition of this tensor, with minimal subspaces
$\Umin{1}{A} = \overline{\mathrm{Im}(A^*)} \subset H$ and
$\Umin{2}{A} = J(\overline{\mathrm{Im}(A)}) \subset H^*$.

\subsection*{Mixed states of rank $r$}

For a mixed state $\rho \in \Drank{r}{H}$ with spectral decomposition
\begin{equation}
\label{eq:qs:mixed_spectral}
    \rho = \sum_{k=1}^r \lambda_k\, v_k \otimes J(v_k),
    \quad \lambda_k > 0, \quad \sum_{k=1}^r \lambda_k = 1,
\end{equation}
the support is $\supp(\rho) = \spann\{v_k : k=1,\ldots,r\} \in \Grass{r}{H}$.
The positive core $\sigma \coloneqq V^*\rho V \in \Dpos{\CC^r}$
(where $V = [v_1|\cdots|v_r]: \CC^r \to H$ is the column matrix of
eigenvectors) is the finite-dimensional density operator on $\CC^r$
with eigenvalues $\lambda_1, \ldots, \lambda_r$.

\begin{proposition}[Rank stratification as Tucker decomposition]
\label{prop:qs:rank_tucker}
The decomposition
\begin{equation}
\label{eq:qs:rank_strat}
    \Dstate{H} = \bigsqcup_{r=1}^\infty \Drank{r}{H}
    \cup \mathcal{D}_\infty(H)
\end{equation}
is the rank stratification of $\Dstate{H}$ by Tucker rank in the
Banach tensor space $(\cTsa{1}(H), \|\cdot\|_1)$. Specifically, viewing
$\rho \in \Drank{r}{H}$ as a tensor in $H \otimes_\pi H^*$ via
$\cT{1}(H) \cong H \otimes_\pi H^*$:
\begin{enumerate}
\item[\textup{(i)}] $r_1(\rho) = r_2(\rho) = r = \mathrm{rank}(\rho)$.
\item[\textup{(ii)}] $\Umin{1}{\rho} = \supp(\rho) \subset H$ and
    $\Umin{2}{\rho} = J(\supp(\rho)) \subset H^*$.
\item[\textup{(iii)}] The trace norm $\|\rho\|_1 = \Tr(\rho) = 1$
    is the projective tensor norm, satisfying \condINJ\ with equality
    at all density operators \textup{(}Remark~\ref{rem:app:trace_projective}\textup{)}.
\end{enumerate}
\end{proposition}

\begin{proof}
Parts~(i)--(ii) follow from the identification $\cT{1}(H) \cong H
\otimes_\pi H^*$ and the definition of Tucker rank
(Definition~\ref{def:bounded_tucker_rank}): $r_\alpha(\rho)$ is the dimension
of $\Umin{\alpha}{\rho}$, which for $\alpha = 1$ is the support in $H$
and for $\alpha = 2$ is the support in $H^*$ (both have dimension $r$
by~\eqref{eq:qs:mixed_spectral}). Part~(iii) is
Remark~\ref{rem:app:trace_projective}.
\end{proof}

\begin{corollary}[Weak lower semicontinuity of quantum rank]
\label{cor:qs:rank_lsc}
In the trace-norm topology, rank is weakly lower semicontinuous on
$\Dstate{H}$: if $\rho_n \in \Drank{r}{H}$ and $\|\rho_n - \rho\|_1
\to 0$, then $\mathrm{rank}(\rho) \leq r$.
\end{corollary}

\begin{proof}
This is Theorem~\ref{thm:weak_closed} applied to the Tucker manifold
$\Drank{r}{H} \subset \cTsa{1}(H)$: the trace norm satisfies \condINJ\
(Remark~\ref{rem:app:trace_projective}), and $\cTsa{1}(H)$ is reflexive
(it is a closed subspace of the Banach space $\cT{1}(H)$, which is
reflexive when $H$ is finite-dimensional, and the statement holds in
infinite dimensions by the same argument with the weak topology on
$\cTsa{1}(H)$; see~\cite{BratteliRobinson1987}).
\end{proof}

%------------------------------------------------------------
\section{Entanglement and minimal subspace dimension}
\label{sec:qs:entanglement}
%------------------------------------------------------------

\subsection*{Separability and Tucker rank one}

For a multipartite system $H = \bigotimes_{j=1}^d H_j$ and a pure
state $\rho_v = v \otimes J(v)$, Theorem~\ref{thm:qs:minimal_partial_trace}
gives a complete characterisation of separability:

\begin{proposition}[Separability and Tucker rank]
\label{prop:qs:separability}
Let $v \in \bigotimes_{j=1}^d H_j$ with $\|v\|=1$ and let $\alpha
\subset D$ be a subset of parties. The following are equivalent:
\begin{enumerate}
\item[\textup{(a)}] $v$ is separable across the bipartition
    $\alpha | \alpha^c$: $v = u \otimes w$ for some $u \in
    \bigotimes_{j \in \alpha} H_j$, $w \in \bigotimes_{j \in \alpha^c} H_j$.
\item[\textup{(b)}] $r_\alpha(v) = 1$.
\item[\textup{(c)}] $\mathrm{rank}(\rho_\alpha) = 1$, i.e., the reduced
    state $\rho_\alpha = \TrS{\alpha^c}(\rho_v)$ is a pure state.
\item[\textup{(d)}] $\Umin{\alpha}{v}$ is one-dimensional.
\end{enumerate}
\end{proposition}

\begin{proof}
(a)$\Leftrightarrow$(b): by definition of Tucker rank
(Definition~\ref{def:bounded_tucker_rank}).
(b)$\Leftrightarrow$(c): Theorem~\ref{thm:qs:minimal_partial_trace}(b).
(b)$\Leftrightarrow$(d): $r_\alpha(v) = \dim\Umin{\alpha}{v}$.
\end{proof}

\subsection*{Entanglement entropy and minimal subspaces}

The \emph{entanglement entropy} of the bipartition $\alpha|\alpha^c$
is the von Neumann entropy of the reduced state:
\begin{equation}
\label{eq:qs:entanglement_entropy}
    S_\alpha(v) \coloneqq -\Tr\bigl(\rho_\alpha \log \rho_\alpha\bigr)
    = -\sum_{k=1}^{r_\alpha(v)} \lambda_k^{(\alpha)}(v)
    \log \lambda_k^{(\alpha)}(v),
\end{equation}
where $\{\lambda_k^{(\alpha)}(v)\}$ are the eigenvalues of $\rho_\alpha$
from Theorem~\ref{thm:qs:minimal_partial_trace}(c). Since
$\lambda_k^{(\alpha)} > 0$ and $\sum_k \lambda_k^{(\alpha)} = 1$,
we have $0 \leq S_\alpha(v) \leq \log r_\alpha(v)$, with:
\begin{itemize}
\item $S_\alpha(v) = 0$ iff $r_\alpha(v) = 1$ (separable, no entanglement).
\item $S_\alpha(v) = \log r_\alpha(v)$ iff all eigenvalues are equal
    $\lambda_k = 1/r_\alpha(v)$ (maximally entangled).
\end{itemize}
The Tucker rank $r_\alpha(v)$ is thus the \emph{Schmidt rank} of the
bipartition, which upper-bounds the entanglement entropy:
$S_\alpha(v) \leq \log r_\alpha(v)$.

\begin{remark}[The Tucker manifold as a variational class for entanglement]
\label{rem:qs:tucker_entanglement}
The Tucker manifold $\Tucker{\fr}{H}$ for $H = \bigotimes_j H_j$
and $\fr = (r_\alpha)_{\alpha \in D}$ is the set of pure states $v$
with prescribed bipartite Schmidt ranks. States on $\Tucker{\fr}{H}$
have entanglement entropies bounded by $S_\alpha \leq \log r_\alpha$
for all bipartitions $\alpha$. The Dirac--Frenkel principle of
Chapter~\ref{chap:DF}, applied on $\Tucker{\fr}{H}$, yields a
variational dynamics that preserves the entanglement structure
$(r_\alpha)_\alpha$ in the sense that the Schmidt ranks remain
equal to $\fr$ along the trajectory.
\end{remark}

%------------------------------------------------------------
\section{The Banach manifold structure of $\Drank{r}{H}$}
\label{sec:qs:manifold}
%------------------------------------------------------------

The geometry of $\Drank{r}{H}$ is a special case of the Tucker manifold
theory of Chapter~\ref{chap:tucker_manifold}, applied to the Banach
tensor space $(\cTsa{1}(H), \|\cdot\|_1) \cong (H \otimes_\pi H^*,
\|\cdot\|_\pi)$.

\subsection*{The atlas}

Fix $\rho_0 \in \Drank{r}{H}$ and let $S_0 = \supp(\rho_0) \in
\Grass{r}{H}$. Choose an isometry $V_0: \CC^r \to H$ with
$\mathrm{Im}(V_0) = S_0$, and set $P_0 = V_0 V_0^*$ (orthogonal
projection onto $S_0$). For $A \in \cL(\CC^r, S_0^\perp)$ with
$\|A\| < 1$, define the isometry
\begin{equation}
\label{eq:qs:VA}
    V_A \coloneqq V_0(I + A^*A)^{-1/2} + \iota A (I + A^*A)^{-1/2}
    \in \cL(\CC^r, H),
    \quad V_A^* V_A = I_{\CC^r},
\end{equation}
where $\iota: S_0^\perp \hookrightarrow H$ is the inclusion. The chart
map
\begin{equation}
\label{eq:qs:chart}
    \Phi_{S_0}: \cL(\CC^r, S_0^\perp) \times \Dpos{\CC^r}
    \longrightarrow \Drank{r}{H},
    \quad (A, \sigma) \longmapsto V_A \sigma V_A^*,
\end{equation}
is the quantum analogue of the Tucker chart map $\chartmap{\bv}$
of Section~\ref{sec:tucker_manifold:charts}: the Grassmann variable
$A \in \cL(\CC^r, S_0^\perp)$ encodes the motion of the support
(analogous to $L_\alpha \in \cL(\Umin{\alpha}{\bv}, \Wmin{\alpha}{\bv})$),
and the core variable $\sigma \in \Dpos{\CC^r}$ encodes the positive
density on the support (analogous to the core tensor $C^{(D)} \in
\RR^{\times r_\alpha}_*$).

\begin{theorem}[Banach manifold structure of $\Drank{r}{H}$,
{\cite[Theorem~3.5]{AcevedoFalco2026}}]
\label{thm:qs:manifold}
For each $r \in \NN$, the set $\Drank{r}{H}$ admits a smooth
$\mathcal{C}^\infty$-Banach manifold structure in the sense of
Lang~\cite{Lang1999}. Specifically:
\begin{enumerate}
\item[\textup{(i)}] For each $\rho_0 \in \Drank{r}{H}$, the map
    $\Phi_{S_0}$ in~\eqref{eq:qs:chart} is a $\mathcal{C}^\infty$
    diffeomorphism from a neighbourhood of $(0, V_0^*\rho_0 V_0)$ in
    $\cL(\CC^r, S_0^\perp) \times \Dpos{\CC^r}$ onto a neighbourhood
    of $\rho_0$ in $\Drank{r}{H}$.
\item[\textup{(ii)}] The transition maps between overlapping charts are
    $\mathcal{C}^\infty$ (analytic over $\CC$).
\item[\textup{(iii)}] The tangent space at $\rho \in \Drank{r}{H}$ is
    \begin{equation}
    \label{eq:qs:tangent}
        T_\rho \Drank{r}{H}
        = \bigl\{ \dot\rho \in \cTsa{1}(H) :
        (I - P)\dot\rho(I - P) = 0,\;
        \Tr(\dot\rho) = 0 \bigr\},
    \end{equation}
    where $P = \PS{\supp(\rho)}$ is the orthogonal projection onto
    $\supp(\rho)$. This is a complemented closed subspace of
    $\cTsa{1}(H)$.
\item[\textup{(iv)}] The inclusion $\iota: \Drank{r}{H} \hookrightarrow
    \cTsa{1}(H)$ is a smooth immersion.
\end{enumerate}
\end{theorem}

\begin{proof}
We prove each part in turn.

\medskip
\noindent\textbf{Part~(i): The chart map is a $\mathcal{C}^\infty$ diffeomorphism.}

Fix $\rho_0 \in \Drank{r}{H}$, $S_0 = \supp(\rho_0)$, and an isometry
$V_0: \CC^r \to H$ with $\mathrm{Im}(V_0) = S_0$.

\emph{Step~1: $\Phi_{S_0}$ is smooth.}
For $\|A\| < 1$, the operator $I_{\CC^r} + A^*A$ is invertible in
$\cL(\CC^r, \CC^r)$ (since $\CC^r$ is finite-dimensional and $A^*A \geq 0$).
The map $A \mapsto (I + A^*A)^{-1/2}$ is real-analytic on operator
norm balls, hence $A \mapsto V_A = (V_0 + \iota A)(I + A^*A)^{-1/2}$
is $\mathcal{C}^\infty$ by composition. Since $\sigma \mapsto V_A
\sigma V_A^*$ is linear in $\sigma$ for fixed $A$, the full map
$(A, \sigma) \mapsto V_A \sigma V_A^*$ is $\mathcal{C}^\infty$.

\emph{Step~2: Values in $\Drank{r}{H}$.}
For any $(A, \sigma)$ in the domain: $V_A \sigma V_A^* \geq 0$ (since
$\sigma > 0$ and $V_A$ is an isometry); $\Tr(V_A \sigma V_A^*) =
\Tr(\sigma V_A^* V_A) = \Tr(\sigma) = 1$ (since $V_A^* V_A =
I_{\CC^r}$); and $\mathrm{rank}(V_A \sigma V_A^*) = r$ (since
$\sigma \succ 0$ and $V_A$ is injective). Hence $\Phi_{S_0}(A,\sigma)
\in \Drank{r}{H}$.

\emph{Step~3: Local bijectivity and smooth inverse.}
At $A = 0$: $V_0 = V_0 (I)^{-1/2} = V_0$ and $\Phi_{S_0}(0,\sigma) =
V_0 \sigma V_0^*$; in particular $\Phi_{S_0}(0, V_0^* \rho_0 V_0) =
V_0 (V_0^* \rho_0 V_0) V_0^* = \rho_0$. For $\rho'$ near $\rho_0$
in $\Drank{r}{H}$, $\supp(\rho')$ is near $S_0$ in $\Grass{r}{H}$,
so there is a unique small $A$ with $\mathrm{Im}(V_A) = \supp(\rho')$.
The inverse is $\Phi_{S_0}^{-1}(\rho') = (A(\rho'), V_A^* \rho' V_A)$,
where $A(\rho') = \chi_{S_0, S_0^\perp}(\supp(\rho'))$ is the
Grassmannian chart map. Since the Grassmannian chart is $\mathcal{C}^\infty$
(Theorem~\ref{thm:grassmann_manifold}) and $\rho' \mapsto V_A^* \rho'
V_A$ is linear in $\rho'$, the inverse is $\mathcal{C}^\infty$.

\medskip
\noindent\textbf{Part~(ii): Smooth transition maps.}

Let $\rho_0, \tilde\rho_0 \in \Drank{r}{H}$ with chart domains
$\mathcal{B}(\rho_0)$ and $\mathcal{B}(\tilde\rho_0)$ overlapping.
The transition map $\Phi_{\tilde S_0}^{-1} \circ \Phi_{S_0}$ is the
composition:
\begin{equation*}
    (A, \sigma) \xmapsto{\Phi_{S_0}} \rho
    \xmapsto{\Phi_{\tilde S_0}^{-1}} (\tilde A, \tilde\sigma),
\end{equation*}
where $\tilde A = \chi_{\tilde S_0, \tilde S_0^\perp}(\supp(\rho))$
is smooth in $\rho$ (Grassmannian transition map, Theorem~\ref{thm:grassmann_manifold}),
and $\tilde\sigma = \tilde V_{\tilde A}^* \rho \tilde V_{\tilde A}$
is smooth in both $\tilde A$ and $\rho$. The composition is therefore
$\mathcal{C}^\infty$.

\medskip
\noindent\textbf{Part~(iii): Tangent space characterisation.}

Let $t \mapsto \rho(t) = \Phi_{S_0}(A(t), \sigma(t))$ be a
$\mathcal{C}^1$ curve with $A(0) = 0$ and $\sigma(0) = V_0^* \rho_0
V_0$. Differentiating $\rho(t) = V_{A(t)} \sigma(t) V_{A(t)}^*$ at
$t=0$:
\begin{equation*}
    \dot\rho(0) = \dot V_0 \sigma_0 V_0^* + V_0 \dot\sigma_0 V_0^*
    + V_0 \sigma_0 \dot V_0^*,
\end{equation*}
where $\dot V_0 = \frac{d}{dt}V_{A(t)}|_{t=0}$. Since $V_{A(t)}^*
V_{A(t)} = I$ implies $\dot V_0^* V_0 + V_0^* \dot V_0 = 0$ (so
$V_0^* \dot V_0$ is skew-Hermitian), and $\mathrm{Im}(V_0) = S_0$
(so $\dot V_0 \in \cL(\CC^r, S_0^\perp)$ by the Grassmannian chart),
one computes:
\begin{equation*}
    (I - P)\dot\rho(0)(I - P)
    = (I-P)\dot V_0 \sigma_0 V_0^*(I-P) = 0,
\end{equation*}
since $V_0^*(I-P) = 0$ (the range of $V_0$ is $S_0$). Also
$\Tr(\dot\rho(0)) = \Tr(\dot\sigma_0) + 2\mathrm{Re}\,\Tr(\sigma_0
V_0^* \dot V_0) = 0$ if $\Tr(\dot\sigma_0) = 0$ and $V_0^*\dot V_0$
is skew-Hermitian. This shows the inclusion
$\dot\rho(0) \in T_\rho \Drank{r}{H} \supset$ the set in~\eqref{eq:qs:tangent}.
The reverse inclusion follows by constructing explicit curves realising
each tangent vector of the given form (Step~3 of the proof
in~\cite{AcevedoFalco2026}).

\medskip
\noindent\textbf{Part~(iv): Smooth immersion.}

The differential of $\iota: \Drank{r}{H} \hookrightarrow \cTsa{1}(H)$
at $\rho$ is the map $d\iota_\rho: T_\rho \Drank{r}{H} \to \cTsa{1}(H)$,
$d\iota_\rho(\dot\rho) = \dot\rho$, which is injective by definition.
The image $T_\rho \Drank{r}{H}$ is a complemented subspace of
$\cTsa{1}(H)$ by the explicit bounded projection $\hat P_\rho$ of
Proposition~\ref{prop:qs:projector} below, satisfying
$\|\hat P_\rho\| \leq 4 < \infty$. Hence $d\iota_\rho$ has closed
split image (Proposition~\ref{prop:split_equiv}), i.e., $\iota$ is
an immersion.
\end{proof}

\subsection*{The tangent projector}

\begin{proposition}[Bounded projector onto $T_\rho \Drank{r}{H}$,
{\cite[Proposition~3.4]{AcevedoFalco2026}}]
\label{prop:qs:projector}
Let $\rho \in \Drank{r}{H}$ and $P = \PS{\supp(\rho)}$. The map
\begin{equation}
\label{eq:qs:projector}
    \hat{P}_\rho: \cTsa{1}(H) \longrightarrow T_\rho \Drank{r}{H},
    \quad
    \hat{P}_\rho(X) \coloneqq PXP + PX(I-P) + (I-P)XP
    - \frac{\Tr(PXP)}{r} P,
\end{equation}
is a bounded linear projection with $\|\hat{P}_\rho\|_{\cL(\cTsa{1}(H))}
\leq 4$, and $\hat{P}_\rho(X) = X$ for all $X \in T_\rho \Drank{r}{H}$.
\end{proposition}

\begin{proof}
Write $X \in \cTsa{1}(H)$ in block form with respect to $H = S \oplus S^\perp$:
\begin{equation*}
    X = \begin{pmatrix} X_{11} & X_{12}^* \\ X_{12} & X_{22}
    \end{pmatrix},
    \quad X_{11} \in \cTsa{1}(S),\; X_{12} \in \cT{1}(S,S^\perp),\;
    X_{22} \in \cTsa{1}(S^\perp).
\end{equation*}
By direct computation using $P^2 = P = P^*$:
\begin{align*}
    PXP &= \begin{pmatrix} X_{11} & 0 \\ 0 & 0 \end{pmatrix}, &
    PX(I-P) &= \begin{pmatrix} 0 & X_{12}^* \\ 0 & 0 \end{pmatrix}, \\
    (I-P)XP &= \begin{pmatrix} 0 & 0 \\ X_{12} & 0 \end{pmatrix}, &
    \frac{\Tr(PXP)}{r}P &= \frac{\Tr(X_{11})}{r}
    \begin{pmatrix} I_S & 0 \\ 0 & 0 \end{pmatrix}.
\end{align*}
Hence
\begin{equation}
\label{eq:qs:projector_block}
    \hat P_\rho(X)
    = \begin{pmatrix} X_{11} - \frac{\Tr(X_{11})}{r}I_S & X_{12}^* \\
    X_{12} & 0 \end{pmatrix}.
\end{equation}
\emph{Idempotency:} Applying $\hat P_\rho$ to~\eqref{eq:qs:projector_block}:
the $(S,S)$ block is $X_{11} - \frac{\Tr(X_{11})}{r}I_S - \frac{1}{r}
\Tr(X_{11} - \frac{\Tr(X_{11})}{r}I_S)I_S = X_{11} - \frac{\Tr(X_{11})}{r}I_S$
(since $\Tr(X_{11} - \frac{\Tr(X_{11})}{r}I_S) = \Tr(X_{11}) - \Tr(X_{11}) = 0$),
and the off-diagonal blocks are unchanged. So $\hat P_\rho^2 = \hat P_\rho$.

\emph{Image:} From~\eqref{eq:qs:projector_block}, $\hat P_\rho(X)$ has
$(S^\perp,S^\perp)$ block zero and $\Tr(\hat P_\rho(X)) = \Tr(X_{11}
- \frac{\Tr(X_{11})}{r}I_S) + 0 = 0$. So $\hat P_\rho(X) \in
T_\rho \Drank{r}{H}$.

\emph{$\hat P_\rho = \mathrm{Id}$ on $T_\rho\Drank{r}{H}$:} If $Y \in
T_\rho\Drank{r}{H}$, its block form is $Y_{22} = 0$ and $\Tr(Y_{11})=0$
(by~\eqref{eq:qs:tangent}). Then $\frac{\Tr(PYP)}{r} = \frac{\Tr(Y_{11})}{r} = 0$,
so $\hat P_\rho(Y) = Y$.

\emph{Norm bound:} Using the ideal property $\|BAC\|_1 \leq \|B\|\|A\|_1\|C\|$:
\begin{align*}
    \|\hat P_\rho(X)\|_1
    &\leq \|PXP\|_1 + \|PX(I-P)\|_1 + \|(I-P)XP\|_1
    + \tfrac{|\Tr(PXP)|}{r}\|P\|_1 \\
    &\leq \|X\|_1 + \|X\|_1 + \|X\|_1
    + \tfrac{\|PXP\|_1}{r} \cdot r \\
    &\leq 4\|X\|_1,
\end{align*}
where we used $|\Tr(PXP)| \leq \|PXP\|_1 \leq \|X\|_1$ and
$\|P\|_1 = \Tr(P) = r$.
\end{proof}

\begin{remark}[Comparison with the Tucker projector]
\label{rem:qs:projector_comparison}
The projector $\hat{P}_\rho$ in~\eqref{eq:qs:projector} is the
quantum analogue of the metric projection $\Pmetric{\bv_r}$ onto
$\ZD{\bv_r}$ from Chapter~\ref{chap:DF}. The explicit formula in
terms of blocks with respect to $H = S \oplus S^\perp$ is the
quantum version of the tangent space formula~\eqref{eq:tangent_formula}:
the $PXP$ block corresponds to the core variation $\dot{C}^{(D)}$,
and the off-diagonal blocks correspond to the subspace velocity
$\dot{u}^{(\alpha)}_{i_\alpha} \in \Wmin{\alpha}{\bv_r}$.
The trace correction $-\frac{\Tr(PXP)}{r}P$ is the additional
constraint $\Tr(\dot\rho) = 0$ that appears because $\Drank{r}{H}$
lies in the affine hyperplane $\{\Tr = 1\}$, not in the full Banach
space.
\end{remark}

\subsection*{The fibre bundle structure}

The support map $\pi: \Drank{r}{H} \to \Grass{r}{H}$,
$\pi(\rho) = \supp(\rho)$, is the quantum analogue of the rank map
$\varrho_\fr$ of Section~\ref{sec:tucker_manifold:cinfty}.

\begin{theorem}[Fibre bundle, {\cite[Theorem~3.8]{AcevedoFalco2026}}]
\label{thm:qs:bundle}
The support map $\pi: \Drank{r}{H} \to \Grass{r}{H}$ is a smooth
fibre bundle with typical fibre $\Dpos{\CC^r}$, a smooth manifold of
real dimension $r^2 - 1$. The structure group is $U(r)$ (the unitary
group), acting on $\Dpos{\CC^r}$ by $U \cdot \sigma = U \sigma U^*$.
\end{theorem}

This is the quantum specialisation of Theorem~\ref{thm:Tucker_Banach_Manifold}(ii)
and Propositions~\ref{prop:tucker_triv}--\ref{prop:tucker_transition_fns}:
the structure group reduces from $\GL{r}{\RR}$ to $U(r)$ because the
core variable $\sigma \in \Dpos{\CC^r}$ transforms by unitary
conjugation (preserving positivity and trace), rather than by
an arbitrary invertible change of basis.

%------------------------------------------------------------
\section{Best approximation and quantum state compression}
\label{sec:qs:approx}
%------------------------------------------------------------

The best-approximation theorem (Theorem~\ref{thm:best_approx_tucker})
applied to the quantum setting gives the mathematical foundation for
\emph{low-rank approximation of quantum states}.

\begin{theorem}[Existence of best low-rank approximation]
\label{thm:qs:best_approx}
Let $\rho \in \Dstate{H}$ and $r \in \NN$. There exists
$\rho_r \in \bigcup_{k=1}^r \Drank{k}{H}$ such that
\begin{equation}
\label{eq:qs:best_approx}
    \|\rho - \rho_r\|_1
    = \inf_{\sigma \in \bigcup_{k=1}^r \Drank{k}{H}} \|\rho - \sigma\|_1.
\end{equation}
\end{theorem}

\begin{proof}
The set $\bigcup_{k=1}^r \Drank{k}{H}$ is weakly closed in
$\cTsa{1}(H)$ by Corollary~\ref{cor:qs:rank_lsc} (weak lower
semicontinuity of rank). The trace-class space $\cTsa{1}(H)$ is
reflexive when $H$ is finite-dimensional; in infinite dimensions,
the result follows from the explicit formula: the best rank-$r$
approximation of $\rho$ in trace norm is given by the \emph{spectral
truncation}
\begin{equation}
\label{eq:qs:spectral_trunc}
    \rho_r \coloneqq \frac{1}{\sum_{k=1}^r \lambda_k}
    \sum_{k=1}^r \lambda_k\, v_k \otimes J(v_k),
\end{equation}
where $\lambda_1 \geq \lambda_2 \geq \cdots$ are the eigenvalues of
$\rho$ in decreasing order and $\{v_k\}$ the corresponding eigenvectors.
The error is $\|\rho - \rho_r\|_1 = 2\sum_{k > r} \lambda_k$
(Ky Fan's theorem).
\end{proof}

\begin{remark}[Spectral truncation as Tucker approximation]
\label{rem:qs:ky_fan}
The spectral truncation~\eqref{eq:qs:spectral_trunc} is the quantum
analogue of the HOSVD (Higher-Order SVD) truncation for tensor formats
(Algorithm~\ref{app:alg:hosvd} in Appendix~\ref{app:algorithms}):
in both cases, the best rank-$r$ approximation in the ambient norm is
obtained by projecting onto the dominant eigenspaces/singular subspaces.
For tensors with $d \geq 3$, the analogue of~\eqref{eq:qs:spectral_trunc}
gives a near-best Tucker approximation (quasi-optimality with
constant $\sqrt{d}$), not an exact best approximation; see
Appendix~\ref{app:alg:hosvd}.
\end{remark}

%------------------------------------------------------------
\section{Historical notes}
\label{sec:qs:history}
%------------------------------------------------------------

\paragraph{Schmidt decomposition and entanglement.}
The Schmidt decomposition for bipartite systems was introduced by
Schmidt~(1907) in the context of integral equations, long before its
quantum-mechanical interpretation. Its role in quantum entanglement
was recognised in the 1990s with the development of quantum information
theory. The generalisation to multipartite systems via minimal subspaces
is the contribution of Falcó and Hackbusch~\cite{FH2012}.

\paragraph{Tucker format and quantum many-body physics.}
The connection between the Tucker format and tree tensor networks (TTN)
was developed in the quantum chemistry and condensed matter literature
under the name MCTDH (Multi-Configuration Time-Dependent Hartree) and
its matrix product state (MPS/DMRG) variants. White's DMRG algorithm
\cite{White1992} is algorithmically the DMRG of
Appendix~\ref{app:algorithms}, implemented for quantum spin chains.
The mathematical equivalence between Tucker formats and tensor network
states was clarified in~\cite{Hackbusch2019}.

\paragraph{Banach manifold structure of $\Drank{r}{H}$.}
The explicit Grassmann-Banach manifold model for fixed-rank density
operators in the trace-norm topology, with the support-core charts and
bounded tangent projectors, was developed in the companion
paper~\cite{AcevedoFalco2026}. Earlier work in the Hilbert--Schmidt
(T$_2$) setting includes~\cite{KochLubich2007}; the move to the
trace-class (T$_1$) setting is motivated by the operational
interpretation of the trace norm as the natural metric for state
distinguishability (Remark~\ref{rem:app:trace_projective}).
